% AX-b literature leaderboard.
% Compile with: pdflatex docs/axb_leaderboard.tex
\documentclass[10pt]{article}
\usepackage[margin=1in]{geometry}
\usepackage[T1]{fontenc}
\usepackage{booktabs}
\usepackage{array}
\usepackage{longtable}
\usepackage{hyperref}
\usepackage{caption}
\usepackage{graphicx}

\title{Natural Language Inference Literature Leaderboards}
\author{GraphNLI-AXb project}
\date{October 2026}

\begin{document}
\maketitle

\section*{Scope and protocol}

This document summarizes published results for Natural Language Inference
(NLI) benchmarks used in the GraphNLI-AXb project. NLI evaluates whether a
hypothesis is entailed by, contradicts, or is neutral with respect to a
premise. The covered resources span English, multilingual and Portuguese
evaluation settings: SICK, XNLI, ASSIN, ASSIN2, InferBR and AX-b.

The tables are organized by dataset because their label spaces, splits and
evaluation protocols are different. SICK, XNLI and InferBR generally use
three-way classification and report accuracy. ASSIN uses the three classes
entailment, paraphrase and neutral, while ASSIN2 uses binary entailment versus
non-entailment. Portuguese studies commonly report accuracy and macro-F1.

Results are ranked only within a table and only when their metric and protocol
are sufficiently aligned. In particular, zero-shot, translate-train,
translate-test and supervised fine-tuning are reported explicitly rather than
being silently combined. Values marked as development, cross-dataset transfer
or a different metric are contextual references and should not be interpreted
as direct leaderboard entries.

\section*{AX-b}

AX-b (Broadcoverage Diagnostics) is a diagnostic NLI set derived from GLUE
Diagnostics and included in SuperGLUE. It focuses on linguistic phenomena
such as negation, quantifiers, conjunction, conditionals, monotonicity and
temporal reasoning. SuperGLUE maps contradiction and neutral to the binary
label \texttt{not\_entailment}, so AX-b is evaluated with Matthews
correlation coefficient (MCC) rather than three-way accuracy
\cite{wang2019superglue}.

The official leaderboard reports MCC multiplied by 100. Thus, a leaderboard
value of 72.9 corresponds to MCC $=0.729$ on the usual $[-1,1]$ scale. The
accuracy table is kept separate because OpenCompass reports a different metric
and protocol.

Table~\ref{tab:literature} collects the strongest publicly documented AX-b
references found in the literature review. Results are deliberately not
merged into a single numerical ranking when the metric or evaluation protocol
changes. In particular, the 2024 PanGu-$\pi$ paper reports accuracy through
OpenCompass \cite{opencompass2023}, whereas the official SuperGLUE diagnostic
uses MCC.

\subsection*{MCC leaderboard}
\small
\begin{longtable}{@{}c c p{3.7cm} p{2.5cm} c p{4.3cm}@{}}
\caption{AX-b literature leaderboard. MCC values are multiplied by 100, as in
the official SuperGLUE leaderboard.}\label{tab:literature}\\
\toprule
\textbf{Rank} & \textbf{Year} & \textbf{System} & \textbf{Total parameters} & \textbf{AX-b MCC} & \textbf{Primary reference or provenance} \\
\midrule
1 & 2019 & SuperGLUE human baseline & -- & 76.6 & Human reference \cite{wang2019superglue} \\
2 & 2021 & PAI Albert & N/D & 75.6 & Official leaderboard submission; primary paper not located \\
3 & 2023 & PaLM 540B & 540B & 72.9 & Total size is documented in the literature \cite{tedeschi2023superhuman} \\
4 & 2022 & ST-MoE-32B & 269B & 72.3 & Total size reported in the SPoT comparison table; sparse model \cite{fedus2022stmoe,vu2022spot} \\
5 & 2021 & T5 + UDG & 11B & 69.1 & Total size reported in the SPoT comparison table \cite{wang2021udg,vu2022spot} \\
6 & 2021 & ERNIE 3.0 & 12B & 68.6 & Total size reported in the SPoT comparison table \cite{sun2021ernie,vu2022spot} \\
7 & 2021 & Turing NLR v5 & 5.4B & 67.8 & Total size reported in the SPoT comparison table \cite{gao2021tnlr,vu2022spot} \\
8 & 2022 & SPoT (Frozen T5 1.1) & 11B & 66.9 & Total size reported by the authors \cite{vu2022spot} \\
9 & 2021 & DeBERTa / TuringNLRv4 & 3.1B & 66.7 & Total size reported in the SPoT comparison table \cite{he2021deberta,vu2022spot} \\
10 & 2020 & T5 & 11B & 65.6 & Total size reported in the SPoT comparison table \cite{raffel2020t5,vu2022spot} \\
11 & 2023 & RoBERTa-iCETS & N/D & 61.8 & Official leaderboard submission; primary paper not located \\
12 & 2022 & RoBERTa-mtl-adv & N/D & 58.5 & Official leaderboard submission; primary paper not located \\
13 & 2021 & NEZHA-Plus & N/D & 58.0 & Official leaderboard submission; primary paper not located \\
14 & 2019 & Facebook RoBERTa & N/D & 57.9 & Official model variant and tuning details are not specified in the leaderboard row \cite{liu2019roberta} \\
15 & 2022 & RoBERTa (ensemble) & N/D & 57.6 & Official leaderboard submission; exact ensemble composition not documented \\
16 & 2020 & Snorkel & N/D & 47.6 & Official leaderboard submission, version 1.9 \\
17 & 2019 & BERT++ & N/D & 38.0 & Original baseline; exact parameterization is not separately reported \cite{wang2019superglue} \\
18 & 2021 & iPET (ALBERT), 32 examples & N/D & 36.2 & Primary paper; exact leaderboard checkpoint parameterization is not separately reported \cite{schick2021ipet} \\
19 & 2021 & FSL++ (ALBERT), 32 examples & N/D & 35.9 & Official leaderboard submission; primary paper not located \\
20 & 2020 & BERT-mtl & N/D & 29.6 & Official leaderboard submission; exact backbone size not documented \\
21 & 2019 & BERT-large & 340M & 23.0 & BERT-Large size documented in the baseline setup \cite{wang2019superglue,devlin2019bert} \\
22 & 2020 & GPT-3 few-shot & 175B & 21.1 & Total size reported in the SPoT comparison table \cite{brown2020gpt3,vu2022spot} \\
23 & 2019 & Most Frequent Class & -- & 0.0 & Original baseline \cite{wang2019superglue} \\
24 & 2021 & ADAPET (ALBERT), few-shot & N/D & -0.4 & Primary paper; exact leaderboard checkpoint parameterization is not separately reported \cite{tam2021adapet} \\
25 & 2023 & Vega v2 & N/D & -0.4 & Primary technical report; exact submitted checkpoint size is not stated in the leaderboard row \cite{wang2022vega} \\
26 & 2021 & INSTALL (ALBERT), few-shot & N/D & -0.4 & Official leaderboard submission; primary paper not located \\
27 & 2021 & WARP (ALBERT-XXL), few-shot & 223M & -0.4 & Total size reported in the SPoT comparison table \cite{vu2022spot} \\
28 & 2019 & CBoW & 15M & -0.4 & Total size reported in the SPoT comparison table \cite{vu2022spot} \\
29 & 2021 & Select-step-by-step & N/D & -0.4 & Official leaderboard submission; primary paper not located \\
\bottomrule
\end{longtable}

\subsection*{Accuracy leaderboard}
\begin{longtable}{@{}c c p{3.7cm} p{2.5cm} c p{4.3cm}@{}}
\caption{AX-b results reported with accuracy rather than MCC. Values are
reported by Tang et al.~\cite{tang2024tinylm} using OpenCompass
\cite{opencompass2023}. These values must not be ranked together with
Table~\ref{tab:literature}.}\label{tab:accuracy}\\
\toprule
\textbf{Rank} & \textbf{Year} & \textbf{System} & \textbf{Total parameters} & \textbf{AX-b accuracy} & \textbf{Primary reference or provenance} \\
\midrule
1 & 2024 & PanGu-$\pi$-1.5B Pro & 1.5B & 67.93 & Model size documented in the primary paper; OpenCompass accuracy, not MCC \cite{tang2024tinylm,opencompass2023} \\
2 & 2024 & Qwen-1.8B & 1.8B & 57.79 & Model size documented in the primary paper; OpenCompass accuracy, not MCC \cite{tang2024tinylm,opencompass2023} \\
3 & 2024 & MobileLLaMA-2.7B & 2.7B & 55.80 & Model size is part of the evaluated model identifier; OpenCompass accuracy, not MCC \cite{opencompass2023} \\
4 & 2024 & RWKV-5-1.5B & 1.5B & 54.05 & Model size is part of the evaluated model identifier; OpenCompass accuracy, not MCC \cite{opencompass2023} \\
5 & 2024 & Open-LLaMA-3B & 3B & 52.72 & Model size is part of the evaluated model identifier; OpenCompass accuracy, not MCC \cite{opencompass2023} \\
6 & 2022 & BLOOM-1B1, one-shot & 1.1B & 52.72 & One-shot evaluation; BigScience artifact \cite{bigscience2022bloom} \\
7 & 2024 & Chinese-LLaMA2-1.3B & 1.3B & 47.46 & Model size is part of the evaluated model identifier; OpenCompass accuracy, not MCC \cite{opencompass2023} \\
8 & 2024 & TinyLLaMA-1.1B & 1.1B & 45.47 & Model size is part of the evaluated model identifier; OpenCompass accuracy, not MCC \cite{opencompass2023} \\
9 & 2024 & MobileLLaMA-1.4B & 1.4B & 45.20 & Model size is part of the evaluated model identifier; OpenCompass accuracy, not MCC \cite{opencompass2023} \\
10 & 2024 & Phi-1.3B & 1.3B & 44.20 & Model size is part of the evaluated model identifier; OpenCompass accuracy, not MCC \cite{opencompass2023} \\
11 & 2024 & Sheared-LLaMA-1.3B & 1.3B & 43.57 & Model size is part of the evaluated model identifier; OpenCompass accuracy, not MCC \cite{opencompass2023} \\
11 & 2024 & Phi-2.7B & 2.7B & 43.57 & Model size is part of the evaluated model identifier; OpenCompass accuracy, not MCC \cite{opencompass2023} \\
13 & 2024 & PanGu-$\pi$-1B & 1B & 43.48 & Model size documented in the primary paper; OpenCompass accuracy, not MCC \cite{tang2024tinylm,opencompass2023} \\
\bottomrule
\end{longtable}

\section*{SICK}
\begin{longtable}{@{}c p{3.2cm} p{3.2cm} p{3.0cm} c p{3.0cm}@{}}
\caption{SICK test-set leaderboard. All systems are trained on the original
SICK training split and evaluated on the original three-way SICK test set.}
\label{tab:sick-literature}\\
\toprule
\textbf{Rank} & \textbf{Year} & \textbf{System} & \textbf{Parameters} & \textbf{Accuracy} & \textbf{Reference} \\
\midrule
1 & 2023 & RoBERTa-Large & 355M & 91.50 & Sharma et al. \cite{sharma2023nli} \\
2 & 2023 & RoBERTa-Large & 355M & 90.5 & Kalouli et al. \cite{kalouli2023sick} \\
3 & 2023 & ALBERT-Large & 17M & 89.85 & Sharma et al. \cite{sharma2023nli} \\
4 & 2023 & BERT-Large & 340M & 89.85 & Sharma et al. \cite{sharma2023nli} \\
5 & 2023 & BERT-Base & 110M & 85.71 & Sharma et al. \cite{sharma2023nli} \\
6 & 2023 & BERT-Large & 340M & 85.3 & Kalouli et al. \cite{kalouli2023sick} \\
7 & 2023 & DAM & 382K & 83.37 & Sharma et al. \cite{sharma2023nli} \\
8 & 2023 & ESIM & 4.3M & 80.81 & Sharma et al. \cite{sharma2023nli} \\
\bottomrule
\end{longtable}

The Sharma et al. results are an independent reimplementation study, while
Kalouli et al. use SICK to analyze annotation regimes. The small numerical
difference between the two RoBERTa/BERT evaluations should therefore not be
interpreted as a strict leaderboard improvement.

\section*{XNLI}
\begin{longtable}{@{}c p{3.1cm} p{2.9cm} p{2.7cm} c p{3.5cm}@{}}
\caption{XNLI average accuracy across the 15 official languages.}
\label{tab:xnli-literature}\\
\toprule
\textbf{Rank} & \textbf{Year} & \textbf{System} & \textbf{Protocol} & \textbf{Accuracy} & \textbf{Reference} \\
\midrule
1 & 2025 & Gemma 2 9B + FLARE & Translate-train & 85.01 & Borchert et al. \cite{borchert2025flare} \\
2 & 2023 & RoBERTa-Large + NLLB & Translate-test & 83.8 & Artetxe et al. \cite{artetxe2023mt} \\
3 & 2022 & XLM-R-Large + PCT & Full-shot + translate-train-all & 83.6 & Qi et al. \cite{qi2022pct} \\
4 & 2020 & XLM-R-Large & Translate-train-all & 82.4 & Conneau et al. \cite{conneau2020xlmr} \\
5 & 2020 & XLM-R-Large & Zero-shot transfer & 79.2 & Xue et al. \cite{xue2021mt5} \\
6 & 2020 & mBERT & Zero-shot transfer & 65.4 & Xue et al. \cite{xue2021mt5} \\
\bottomrule
\end{longtable}

The XNLI table intentionally exposes the protocol. Translate-train and
translate-test use additional machine-translated inputs, whereas zero-shot
transfer uses English training data and evaluates directly in the target
languages. These regimes should be compared only within the same protocol.

\section*{ASSIN}
\begin{longtable}{@{}c p{3.2cm} p{3.4cm} c c c p{3.2cm}@{}}
\caption{ASSIN three-class NLI leaderboard on the combined Portuguese
European/Brazilian test setting. The classes are entailment, paraphrase and
neutral.}
\label{tab:assin-literature}\\
\toprule
\textbf{Rank} & \textbf{Year} & \textbf{System} & \textbf{Parameters} & \textbf{Accuracy} & \textbf{Macro-F1} & \textbf{Reference} \\
\midrule
1 & 2020 & ptBERT-Large fine-tuned & 335M & 91.0 & 82.0 & Fialho et al. \cite{fialho2020} \\
2 & 2020 & ptBERT-Base fine-tuned & 110M & 90.0 & 81.0 & Fialho et al. \cite{fialho2020} \\
3 & 2020 & mBERT-Base fine-tuned & 178M & 90.0 & 80.0 & Fialho et al. \cite{fialho2020} \\
4 & 2016 & ASSIN challenge best system & N/D & 83.85 & 71.0 & Fonseca et al. \cite{fonseca2016assin} \\
\bottomrule
\end{longtable}

ASSIN is not identical to the binary NLI formulation used by ASSIN2 or
InferBR. In particular, paraphrase is a separate class here; it should not be
collapsed into neutral when comparing these values.

\section*{ASSIN2}
\begin{longtable}{@{}c p{3.1cm} p{3.4cm} c c c p{3.0cm}@{}}
\caption{ASSIN2 binary NLI leaderboard. The task labels are entailment and
non-entailment; the official test split is used.}
\label{tab:assin2-literature}\\
\toprule
\textbf{Rank} & \textbf{Year} & \textbf{System} & \textbf{Parameters} & \textbf{Accuracy} & \textbf{Macro-F1} & \textbf{Reference} \\
\midrule
1 & 2020 & BERTimbau-Large fine-tuned & 335M & 90.0 & 90.0 & Souza et al. \cite{souza2020bertimbau} \\
2 & 2026 & BERTomelo-Large fine-tuned & N/D & 89.26 & 89.21 & Oliveira et al. \cite{oliveira2026bertomelo} \\
3 & 2020 & BERTimbau-Base fine-tuned & 110M & 89.2 & 89.2 & Souza et al. \cite{souza2020bertimbau} \\
4 & 2020 & mBERT fine-tuned & 178M & 86.8 & 86.8 & Souza et al. \cite{souza2020bertimbau} \\
5 & 2020 & mBERT + RoBERTa-Large ensemble & N/D & 88.3 & 88.3 & Rodrigues et al. \cite{rodrigues2020assin2} \\
\bottomrule
\end{longtable}

NorBERTo reports approximately 0.904 entailment F1 on ASSIN2, but this is a
per-class F1 rather than macro-F1. It is therefore recorded as a recent
reference, not inserted into the ranking above \cite{silva2026norberto}.

\section*{InferBR}
\begin{longtable}{@{}c p{3.2cm} p{3.5cm} c c c p{3.0cm}@{}}
\caption{InferBR three-class NLI leaderboard. Results are from the original
InferBR evaluation protocol, with BERTimbau fine-tuned on each dataset's
training and validation splits.}
\label{tab:inferbr-literature}\\
\toprule
\textbf{Rank} & \textbf{Year} & \textbf{System} & \textbf{Parameters} & \textbf{Accuracy} & \textbf{Macro-F1} & \textbf{Reference} \\
\midrule
1 & 2024 & BERTimbau fine-tuned on InferBR & N/D & 90.0 & 90.0 & Bencke et al. \cite{bencke2024inferbr} \\
2 & 2024 & BERTimbau trained on SICK-BR, tested on InferBR & N/D & 60.0 & 58.0 & Bencke et al. \cite{bencke2024inferbr} \\
3 & 2024 & BERTimbau trained on InferBR, tested on SICK-BR & N/D & 64.0 & 63.0 & Bencke et al. \cite{bencke2024inferbr} \\
\bottomrule
\end{longtable}

The cross-dataset rows are included as transfer references, not as direct
InferBR leaderboard entries. They illustrate the domain shift between
synthetically generated InferBR examples and SICK-BR examples.

Here ``Total parameters'' refers to the complete model size. ``N/D'' means
that the primary source or official submission does not document the total
parameter count for that exact entry; it is intentionally not inferred from
the model name. For sparse models, total parameters include all experts. The leaderboard is a historical snapshot and mixes official submissions,
human references, few-shot systems, and baselines. The 2024--2026 review found
no newer peer-reviewed result with a clearly documented official AX-b MCC above
72.9; recent papers often omit AX-b or report a different metric and split.

\bibliographystyle{plain}
\begin{thebibliography}{9}

\bibitem{wang2019superglue}
Alex Wang, Yada Pruksachatkun, Nikita Nangia, Amanpreet Singh, Julian
Michael, Felix Hill, Omer Levy, and Samuel R. Bowman.
\newblock SuperGLUE: A stickier benchmark for general-purpose language
understanding systems.
\newblock In \emph{NeurIPS}, 2019. \url{https://arxiv.org/abs/1905.00537}.

\bibitem{devlin2019bert}
Jacob Devlin, Ming-Wei Chang, Kenton Lee, and Kristina Toutanova.
\newblock BERT: Pre-training of deep bidirectional transformers for language understanding.
\newblock In \emph{NAACL}, 2019. \url{https://aclanthology.org/N19-1423/}.

\bibitem{tedeschi2023superhuman}
Simone Tedeschi, Johan Bos, Thierry Declerck, Jan Hajic, Daniel Hershcovich,
Eduard H. Hovy, Alexander Koller, Simon Krek, Steven Schockaert, Rico Sennrich,
Ekaterina Shutova, and Roberto Navigli.
\newblock What's the meaning of superhuman performance in today's NLU?
\newblock In \emph{ACL}, 2023. \url{https://aclanthology.org/2023.acl-long.697/}.

\bibitem{fedus2022stmoe}
William Fedus et al.
\newblock ST-MoE: Designing stable and transferable sparse expert models.
\newblock \emph{arXiv preprint arXiv:2202.08906}, 2022.
\url{https://arxiv.org/abs/2202.08906}.

\bibitem{wang2021udg}
Zirui Wang, Adams Wei Yu, Orhan Firat, and Yuan Cao.
\newblock Towards zero-label language learning.
\newblock \emph{arXiv preprint arXiv:2109.09193}, 2021.
\url{https://arxiv.org/abs/2109.09193}.

\bibitem{sun2021ernie}
Yu Sun et al.
\newblock ERNIE 3.0: Large-scale knowledge enhanced pre-training for language understanding and generation.
\newblock \emph{arXiv preprint arXiv:2107.02137}, 2021.
\url{https://arxiv.org/abs/2107.02137}.

\bibitem{gao2021tnlr}
Jianfeng Gao and Saurabh Tiwary.
\newblock Efficiently and effectively scaling up language model pretraining for best language representation model on GLUE and SuperGLUE.
\newblock Microsoft Research technical report, 2021.
\url{https://www.microsoft.com/en-us/research/blog/efficiently-and-effectively-scaling-up-language-model-pretraining-for-best-language-representation-model-on-glue-and-superglue/}.

\bibitem{vu2022spot}
Tu Vu et al.
\newblock SPoT: Better frozen model adaptation through soft prompt transfer.
\newblock In \emph{ACL}, 2022.
\url{https://aclanthology.org/2022.acl-long.346/}.

\bibitem{he2021deberta}
Pengcheng He et al.
\newblock DeBERTa: Decoding-enhanced BERT with disentangled attention.
\newblock In \emph{ICLR}, 2021.
\url{https://openreview.net/forum?id=XPZIaotutsD}.

\bibitem{raffel2020t5}
Colin Raffel et al.
\newblock Exploring the limits of transfer learning with a unified text-to-text transformer.
\newblock \emph{Journal of Machine Learning Research}, 2020.
\url{https://jmlr.org/papers/v21/20-074.html}.

\bibitem{liu2019roberta}
Yinhan Liu et al.
\newblock RoBERTa: A robustly optimized BERT pretraining approach.
\newblock \emph{arXiv preprint arXiv:1907.11692}, 2019.
\url{https://arxiv.org/abs/1907.11692}.

\bibitem{schick2021ipet}
Timo Schick and Hinrich Schütze.
\newblock Exploiting cloze questions for few-shot text classification and natural language inference.
\newblock In \emph{EACL}, 2021.
\url{https://aclanthology.org/2021.eacl-main.20/}.

\bibitem{brown2020gpt3}
Tom B. Brown et al.
\newblock Language models are few-shot learners.
\newblock In \emph{NeurIPS}, 2020.
\url{https://papers.nips.cc/paper/2020/hash/1457c0d6bfcb4967418bfb8ac142f64a-Abstract.html}.

\bibitem{tam2021adapet}
David Tam et al.
\newblock Improving and simplifying pattern exploiting training.
\newblock In \emph{EMNLP}, 2021.
\url{https://aclanthology.org/2021.emnlp-main.407/}.

\bibitem{wang2022vega}
Zirui Wang et al.
\newblock Toward efficient language model pretraining and downstream adaptation via self-evolution: A case study on SuperGLUE.
\newblock \emph{arXiv preprint arXiv:2212.01853}, 2022.
\url{https://arxiv.org/abs/2212.01853}.

\bibitem{laskar2023chatgpt}
Md Tahmid Rahman Laskar, M. Saiful Bari, Mizanur Rahman, Md Amran Hossen
Bhuiyan, Shafiq Joty, and Jimmy Xiangji Huang.
\newblock A systematic study and comprehensive evaluation of ChatGPT on
benchmark datasets.
\newblock In \emph{Findings of ACL}, 2023.
\url{https://aclanthology.org/2023.findings-acl.29/}.

\bibitem{tang2024tinylm}
Zhenyu Tang et al.
\newblock Rethinking optimization and architecture for tiny language models.
\newblock In \emph{ICML}, 2024.
\url{https://openreview.net/forum?id=mHIEOZtDDF}.

\bibitem{shi2024dept}
Zhengxiang Shi and Aldo Lipani.
\newblock DEPT: Decomposed prompt tuning for parameter-efficient fine-tuning.
\newblock In \emph{ICLR}, 2024.
\url{https://proceedings.iclr.cc/paper_files/paper/2024/file/e352b765e625934ce86919995e2371aa-Paper-Conference.pdf}.

\bibitem{opencompass2023}
OpenCompass Contributors.
\newblock OpenCompass: A universal evaluation platform for foundation models.
\newblock Technical report, 2023.
\url{https://github.com/open-compass/opencompass}.

\bibitem{bigscience2022bloom}
BigScience Workshop.
\newblock BLOOM: A 176B-parameter open-access multilingual language model.
\newblock \emph{arXiv preprint arXiv:2211.05100}, 2022.
\url{https://arxiv.org/abs/2211.05100}.
The associated one-shot evaluation artifacts are available at
\url{https://huggingface.co/datasets/bigscience/evaluation-results/discussions/5/files}.

\bibitem{anempirical2025merging}
An empirical survey of model merging algorithms.
\newblock In \emph{PACLIC}, 2025.
\url{https://aclanthology.org/anthology-files/pdf/paclic/2025.paclic-1.24.pdf}.

\bibitem{marelli2014sick}
Stefano Marelli et al.
\newblock A SICK cure for the evaluation of compositional distributional semantic models.
\newblock In \emph{LREC}, 2014.
\url{https://aclanthology.org/L14-1314/}.

\bibitem{kalouli2023sick}
Aikaterini-Lida Kalouli, Hai Hu, Alexander F. Webb, Lawrence S. Moss, and
Valeria de Paiva.
\newblock Curing the SICK and Other NLI Maladies.
\newblock \emph{Computational Linguistics}, 49(1), 2023.
\url{https://aclanthology.org/2023.cl-1.5/}.

\bibitem{sharma2023nli}
Saurabh Sharma et al.
\newblock Evaluating deep learning techniques for natural language inference.
\newblock \emph{Applied Sciences}, 13(4), 2023.
\url{https://www.mdpi.com/2076-3417/13/4/2577}.

\bibitem{conneau2018xnli}
Alexis Conneau et al.
\newblock XNLI: Evaluating cross-lingual sentence representations.
\newblock In \emph{EMNLP}, 2018.
\url{https://aclanthology.org/D18-1269/}.

\bibitem{conneau2020xlmr}
Alexis Conneau et al.
\newblock Unsupervised cross-lingual representation learning at scale.
\newblock In \emph{ACL}, 2020.
\url{https://aclanthology.org/2020.acl-main.747/}.

\bibitem{qi2022pct}
Kunxun Qi, Hai Wan, Jianfeng Du, and Haolan Chen.
\newblock Enhancing cross-lingual natural language inference by prompt-learning
from cross-lingual templates.
\newblock In \emph{ACL}, 2022.
\url{https://aclanthology.org/2022.acl-long.134/}.

\bibitem{artetxe2023mt}
Mikel Artetxe, Vedanuj Goswami, Shruti Bhosale, Angela Fan, and Luke
Zettlemoyer.
\newblock Revisiting machine translation for cross-lingual classification.
\newblock In \emph{EMNLP}, 2023.
\url{https://aclanthology.org/2023.emnlp-main.399/}.

\bibitem{borchert2025flare}
Philipp Borchert, Ivan Vuli\'c, Marie-Francine Moens, and Jochen De Weerdt.
\newblock Language fusion for parameter-efficient cross-lingual transfer.
\newblock In \emph{ICLR}, 2025.
\url{https://openreview.net/forum?id=2e6cb95570c91a50feba8b5e6469b6835de7e576}.

\bibitem{xue2021mt5}
Linting Xue et al.
\newblock mT5: A massively multilingual pre-trained text-to-text transformer.
\newblock In \emph{NAACL}, 2021.
\url{https://aclanthology.org/2021.naacl-main.41/}.

\bibitem{fialho2020}
Pedro Fialho, Lu\'isa Coheur, and Paulo Quaresma.
\newblock Benchmarking natural language inference and semantic textual similarity for Portuguese.
\newblock \emph{Information}, 11(10), 2020.
\url{https://doi.org/10.3390/info11100484}.

\bibitem{fonseca2016assin}
Erick Rocha Fonseca, Leandro Borges dos Santos, Marcelo Criscuolo, and
Sandra Maria Alu\'isio.
\newblock Overview of the evaluation of semantic similarity and textual inference.
\newblock \emph{Linguam\'atica}, 8(2), 2016.
\url{https://www.linguamatica.com/index.php/linguamatica/article/view/v8n2-1}.

\bibitem{souza2020bertimbau}
F\'abio Souza, Rodrigo Nogueira, and Roberto Lotufo.
\newblock BERTimbau: Pretrained BERT models for Brazilian Portuguese.
\newblock In \emph{BRACIS}, 2020.
\url{https://doi.org/10.1007/978-3-030-61377-8_28}.

\bibitem{rodrigues2020assin2}
Rodrigues et al.
\newblock ASSIN2 shared-task systems for Portuguese textual entailment.
\newblock In \emph{ASSIN2 shared task}, 2020.
\url{https://sites.google.com/view/assin2}.

\bibitem{oliveira2026bertomelo}
Renn\^e Ruan Alves Oliveira, Gustavo Cordeiro Galv\~ao Van Erven, and Lu\'is
Paulo Faina Garcia.
\newblock BERTomelo: Your Portuguese encoder best friend.
\newblock \emph{arXiv preprint arXiv:2606.28999}, 2026.
\url{https://arxiv.org/abs/2606.28999}.

\bibitem{silva2026norberto}
Enzo S. N. Silva et al.
\newblock NorBERTo: A ModernBERT model trained for Portuguese with 331 billion tokens corpus.
\newblock \emph{arXiv preprint arXiv:2605.00086}, 2026.
\url{https://arxiv.org/abs/2605.00086}.

\bibitem{bencke2024inferbr}
Luciana Bencke, Francielle V. Pereira, Moniele K. Santos, and Viviane P. Moreira.
\newblock InferBR: A natural language inference dataset in Portuguese.
\newblock In \emph{LREC-COLING}, 2024.
\url{https://aclanthology.org/2024.lrec-main.793/}.

\end{thebibliography}

\end{document}
